\documentclass{article}

\usepackage{amssymb,amsmath,amsthm,mathtools}
\usepackage[margin=1in]{geometry}
\usepackage{enumitem}


\newtheorem{theorem}{Theorem}[section]
\newtheorem{lemma}[theorem]{Lemma}

\begin{document}

\noindent
\begin{minipage}[t]{0.68\textwidth}
    {\large 6.1200 \textit{Mathematics for Computer Science}}\\[4pt]
    {\large Problem Set 2}\\[4pt]
    {\large Solutions by Harsh Bhalala}
\end{minipage}%
\hfill
\begin{minipage}[t]{0.27\textwidth}
    \raggedleft
    {\large 1, September 2026}
\end{minipage}

\vspace{6pt}

\noindent\rule{\textwidth}{0.5pt}

\section*{Problem 1. Outlining Proofs}
\noindent\textbf{(a)} \textbf{Theorem:} “Every natural number is tall or wide, but not both.”

\begin{proof}[Proof by cases.]
Suppose $n$ is any natural number. We must show that it cannot be tall and wide at the same time.
To show n must be either tall or wide assume n is not tall [TODO: prove that n is wide].
now all that remain is to show that n cannot be both tall and wide by caces.

\textbf{Case I:} $n$ is tall. 
[TODO: show that $n$ cannot be wide.]

\textbf{Case II:} $n$ is wide.
[TODO: show that $n$ cannot be tall.]

Thus, by cases, we have shown that the theorem holds for every natural numbers.
\end{proof}

\noindent\textbf{(b)} \textbf{Theorem:} “There exists a happy natural number between 10 and 1000 that has no happy proper divisors.”

\begin{proof}[Proof by example.]
Theorem holds for n = 25,

Show that 25 is happy natural number between 10 and 1000 because [TODO].
    
Also prove that 25 has no happy proper divisors [TODO].

\end{proof}

\noindent\textbf{(c)} \textbf{Theorem:} “There are no devious natural numbers.”

\[P(n) \coloneqq \text{$n$ is not devious}\]

\begin{proof}[Proof by induction.]
    
    \textbf{Induction hypothesis.}$P(n)$ holds for every $n \in \mathbb{N}$

    \textbf{Base Cases.}
    
    \begin{itemize}[leftmargin=\parindent]
        \item $P(0)$ holds because [TODO].
        \item $P(1)$ holds because [TODO].
    \end{itemize}
    
    \textbf{Inductive step:} Assume $P(n-1)$ holds for some arbitrary $n \in \mathbb{N}$,
    now to prove that $p(n)$ also holds because 
    
    [TODO].

    By induction principle, for every natural number P(n) holds.
\end{proof}

\section*{Problem 2. Spotting Errors}

\noindent\textbf{(a)} Brynmor makes mistake while Plugging equations into equality. he fotgot to flip the inequalitie when multiplying by $-8$,
the correct resultant equation would be $6x_i + 8x_{i-1} \geq 6 \cdot 2^{i-1} + 8 \cdot 2^{i-1}$ which is not sufficient to 
prove the inductive step.

\vspace{\baselineskip}
\noindent\textbf{(b)} In given proof during inductive steps Erik stated that $x_i > 0$ which was yet to be proven.

\vspace{\baselineskip}
\noindent\textbf{(c)} Zach's inductive step fails on $i = 1$ because queation used in inductive step $x_{i+1} = 6x_i - 8x_{i-1}$ would 
would requre assumption that $S(0 )$ to use value $x_8$ which is not defined.

\pagebreak

\vspace{\baselineskip}
\noindent\textbf{(d)}


\begin{proof}
    We will use induction on stronger hypothesis $R(i) := \text{``} x_i \geq 3x{i-1}\text{''}$ AND $x_i > 0$ ,proving that 
    $R(i)$ holds for $n \geq 2$.
    \newline\emph{Base Case:} $R(2) = 13 > 12$, and $13 > 0$, base case holds.
    \newline\emph{Inductive step:} Suppose for $n \geq 2$ R(i) holds and let's show that R(i+1) holds.
    Consider $x_{x+1} = 6x_i - 8x_{i-1}$ also from induction we know that $x_i \geq 3x_{i-1}$, So $ \frac{8}{3}x \geq 8x_{i-1}$ ,now plug that into 
    $x_{i+1} = 6x_i - 8x_{i-1} = \frac{10}{3}x_i + \frac{8}{3}x_i - 8x_{i-1}$ and you get $x_{i+1} \geq \frac{10}{3}x_i$. also $x_i > 0$ which means $\frac{10}{3}x_i > 3x_i$, $x_{i+1} \geq 3x_i$ and here $x_i > 0$ so, $x_{i+1} > 0$.
    That completes the Inductive step,so $R(i)$ is true for all $i \geq 2$.
    
    \noindent$R(i)$ is true for $i \geq 2$ and $x_1 = 4$,so for all $i \geq 1$, $x_i > 0$ as desire.
\end{proof}

\section*{Problem 3. Integer Multiplication from Bit Shifts}

\noindent\textbf{(a)} start state to final state:
\begin{align*}
    Q_1(5,22,0) &\longrightarrow Q_2(10,11,0) \longrightarrow Q_3(20,5,10) \longrightarrow Q_4(40,2,30) \longrightarrow Q_5(80,1,30) \longrightarrow Q_f(160,0,110)
\end{align*}

in final state $a_f = 110$ as desired, hence state machine computes $5 \cdot 22$ correctly.

\vspace{\baselineskip}
\noindent\textbf{(b)} Theorem:- $P(r, s, a) := \text{``}rs + a= xy\text{''}$ is invariant.
\begin{proof}[Proof by invariant principle]:
    For start state $Q_0$ where r = x,s = y,a = 0 $P(x,y,0)$ holds.
    Now to prove $P$ is preserved we assume $P(r,s,a)$ holds (i.e: $rs + a = xy$) and that $(r,s,a) \rightarrow (r_t,s_t,a_t)$.We must prove that $P((r_t,s_t,a_t))$ 
    holds, that is,
    \[r_ts_t + a_t = xy\]
    we also  prove by cases on transition from $(r,s,a) \rightarrow (r_t,s_t,a_t)$:
    \begin{itemize}
        \item\textbf{Case s is even:} 
            \[(r_t,s_t,a_t) = (2r,s/2,a)\]
            \[2r*s/2 + a = rs + a = xy\]
            which is true by our assumption,so $P((r_t,s_t,a_t))$ holds.
        \item\textbf{Case s is odd:}
            \[(r_t,s_t,a_t) = (2r,(s-1)/2,a+r)\] 
            \[2r*(s-1)/2 + a + r = rs + a = xy\]
            is true rue by our assumption, so $P((r_t,s_t,a_t))$ holds.
    \end{itemize}
    That covers all the transitions and $P$ is preserved, so by invariant principle $P(r,s,a)$ is invariant.
    
    \vspace{\baselineskip}
    \noindent\textbf{(c)} state $(r_f,s_f,a_f)$ is final when $s_f = 0$ 
    
    \vspace{\baselineskip}
    \noindent\textbf{(d)} From part (b) for the final state $P((r_f,s_f,a_f)) = r_f  s_f + a_f = xy$, also using (c) in the expression, we conclude that in any reachable final state $a_f = xy$. 

    \vspace{\baselineskip}
    \noindent\textbf{(e)} $bitcount(s) = \lceil \log_2(s+1)  \rceil$  is derived variable so that $s \in q = (r,s,a)$ where $q \in Q$.

    Now for some arbitrary transition $q \rightarrow q'$, by definition of transition,
    value of $s$ is at least getting halved and so $s_{new}$ is at most $s/2$ and we also that know halving the value
    of a positive integer would lead to its bit count being reduced by 1, meaning $bitcount(s) > bitcount(s_{new})$

    By the given facts, we can conclude that $bitcount(s)$ is strictly decreasing and its output is strictly closed over natural numbers.
    
    This implies that the algorithm will eventually terminate because $bitcout(s)$ is strictly decreasing by each transition and it has a non-negative lower bound, so it must terminate.

    \vspace{\baselineskip}
    \noindent\textbf{(f)} According to Termination Theorem, since $bitcount(s)$ is strictly decreasig derived variable definied over
    set of natual numbers (By Part(e)) it will end in at most $\lceil \log_2(y + 1) \rceil $steps when starting state is $(x,y,0)$.
\end{proof}

\setcounter{section}{4}
\section*{Problem 4. Invariant Enlightenment}

\textbf{Definition:} For our proof, perimeter is the total number of edges between enlightened and unenlightened students.  

\begin{lemma}
    $P(n)$:= With $k$ initially enlightened students, after n days the perimeter will be at most $4k$.
\end{lemma}

\begin{proof}[Proof by induction]:
    
    \textbf{Base Case $P(0)$:} On day 0, if the $k$ enlightened students are seated far apart, their perimeter will be at most $4k$, 
     base case holds.
    
    \textbf{Inductive step:} Assume $P(n)$ is true and now want to show that $P(n+1)$ is also true. 
    On day $n+1$, new students can become enlightened if 2 or more students adjacent to them are enlightened, so when a new student is enlightened, we observe Net change in perimeter by cases on the number of adjacent students.
    \begin{itemize}
        \item \textbf{Case I (2 adjacent):} increase and decrease by 2. 
        \item \textbf{Case II (3 adjacent):} decrease 3 and increase by 1.
        \item \textbf{Case III (4 adjacent):} only decrease by 4.
    \end{itemize}

    All cases by which the perimeter may change are exhausted, and in all of the cases we showed that the total perimeter will stay $4k$ 
    at most, implying $P(n+1)$ and By induction hypothesis $P(n)$.
\end{proof}

\begin{theorem}
If fewer than $n$ students among those in an $n \times n$ arrangment are initially enlightened, and 
additional students reach enlightenment only by being taught by two or more neighbors, then there will be
at least one student who never reaches enlightenment.
\end{theorem}

\begin{proof}
    Let's say that k students were initially enlightened in some $n \times n$ grid, by the end of all possible
    transitions, every student in the grid is enlightened. From Lemma 4.1, we know that the perimeter of such a state can 
    be $4k$ at most, we also know that perimeter of $n \times n$ grid is $4n$.
    \[4k \geq 4n \implies k \geq n\]

    Which means that to completely enlightened grid of $n \times n$ students initially we need at least $n$ students. Which proves the theorem.

\end{proof}

\end{document}
